\paragraph{Explicit fixed-routing adjoint.}
Generic reverse-mode differentiation through the node-loading scan retains more
state than is necessary when link probabilities are fixed. Let $c_\ell$ be the
cotangent of link flow and let $a_i$ be the node-flow adjoint. A reverse
topological pass evaluates
\[
a_i=\sum_{\ell\in\delta^+(i)}P_\ell
\left(c_\ell+a_{h(\ell)}\right).
\]
The derivative with respect to initial flow injected at node $i$ is $a_i$; the
existing OD injection operation then gathers this adjoint at origin nodes. The
custom vector--Jacobian product therefore retains a node vector and the supplied
link cotangent rather than the complete forward scan tape. It deliberately
defines no derivative with respect to routing probabilities and is exposed only
for the fixed-routing path.

On the representative 465-cell full-network destination, ordinary reverse mode
and the explicit adjoint required respectively 12.09 and 6.69~s for a warm
value--gradient evaluation. Their resident memory immediately after that
evaluation was 16.21 and 13.91~GiB. Objective, prediction, and gradient
diagnostics agreed within float32 tolerances. The explicit adjoint is retained
as the fixed-routing differentiation path and as a reference for validating
future sparse-operator construction.

